% Source: https://artofproblemsolving.com/community/c3h3125578_omt_individual_round_p5
% Stable identity: aops-topic-3125578-q01
Suppose the second player had a winning strategy. The first player makes any opening move and thereafter follows that second-player strategy against the opponent, pretending the extra opening mark is absent. Whenever the strategy requests an already owned square, count that existing mark as the prescribed move and place the actual move on another empty square, treating the new mark as extra. An extra own mark cannot obstruct a winning line. Thus the first player would obtain a winning line no later than the supposed second-player strategy dictates, a contradiction. A finite perfect-information game has a winning or drawing strategy for one side; since the second player cannot force a win, the first can guarantee at least a draw.
\begin{asy}
usepackage("amsmath");
size(7cm);
for(int k=0;k<=3;++k){draw((k,0)--(k,3));draw((0,k)--(3,k));}label("$n\times n$ board (example $n=3$)",(1.5,-.5));label("$X$",(.5,.5));
\end{asy}
